\section{Introduction}\label{sec:intro}
The deployment of Large Language Models (LLMs) such as GPT-4 and Gemini in clinical settings has demonstrated their potential to assist clinicians with diagnostic suggestions, literature summaries, and treatment plan generation. Despite these capabilities, generative models are fundamentally statistical next-token predictors. Consequently, they suffer from the ``hallucination'' phenomenon---generating logical-sounding but factual fabrications. In high-stakes clinical settings, an unsubstantiated treatment or diagnostic suggestion can lead to severe adverse patient outcomes.

\subsection{The Arbovirus Differential Diagnosis Dilemma}
To evaluate X-CDS in a highly challenging and clinically realistic scenario, we focus on emerging arboviruses: \textbf{Zika, Dengue, Chikungunya, and West Nile Virus}. Differentiating these pathogens represents a classic diagnostic challenge in tropical medicine due to their overlapping acute presentations (fever, rash, and joint pain). 

However, misdiagnosis carries extreme clinical risk:
\begin{itemize}
    \item A patient presenting with severe joint pain may be diagnosed with Chikungunya and prescribed non-steroidal anti-inflammatory drugs (NSAIDs) like Aspirin or Ibuprofen for pain relief.
    \item If the patient actually has Dengue, administering NSAIDs can impair platelet function and trigger severe, potentially fatal internal hemorrhages (Dengue Hemorrhagic Fever).
    \item Similarly, misdiagnosing Zika in pregnant patients can lead to missed screening opportunities for microcephaly and Congenital Zika Syndrome (CZS).
\end{itemize}

This environment perfectly simulates the high-stakes clinical scenarios where LLM hallucinations or slight inaccuracies are unacceptable. Traditional Retrieval-Augmented Generation (RAG) mitigates this by injecting relevant medical literature into the model's prompt context \citep{lewis2020retrieval}. However, standard RAG still fails if:
\begin{enumerate}
    \item The retrieval pipeline fails to capture critical clinical facts (poor recall).
    \item The generator ignores the injected context or synthesizes claims that lack an explicit source citation (hallucination).
\end{enumerate}

To address these vulnerabilities, we present the \textbf{X-CDS} framework. The primary contributions of this paper are:
\begin{enumerate}
    \item A dual-channel \textbf{Hybrid Retrieval} pipeline combining dense vector embeddings and sparse keyword indices merged via Reciprocal Rank Fusion (RRF) and optimized using a transformer-based Cross-Encoder re-ranker.
    \item A stateful \textbf{LangGraph Self-Correction} workflow that programmatically checks citation alignment using character-level token-overlap algorithms.
    \item An empirical evaluation of the pipeline's effectiveness using the Ragas metric framework, showing that programmatic self-correction loops yield directional improvements in clinical safety.
\end{enumerate}
